% chapters/04/section-06.tex
\section{Yang Belum Ada pada Siklus Ini}
\label{sec:belum-ada}

Siklus permainan sudah berjalan, tetapi permainannya belum. Bagian penutup
ini menjelaskan apa yang masih kurang dan bagaimana kekurangan itu dilengkapi
pada bab berikutnya.

\subsection{Siklus Berjalan, tetapi Tidak Ada yang Bergerak}

Program pada bab ini memperbarui keadaan berupa penghitung bingkai, dan
penayangannya hanya menampilkan angka serta sebuah batang penanda. Tidak ada
objek yang berpindah tempat, karena perpindahan objek memerlukan dua hal yang
belum dibahas.

\begin{enumerate}
  \item \textbf{Keadaan yang lebih kaya.} Posisi dan kecepatan setiap objek
        harus disimpan dan diperbarui setiap bingkai.
  \item \textbf{Penggunaan selang waktu untuk perhitungan.} Program ini
        sudah mengukur selang waktu melalui \code{clock.tick()}, tetapi nilai
        itu baru dipakai untuk menampilkan waktu berjalan. Pada Bab~V nilai
        yang sama dipakai sebagai pengali kecepatan, sehingga objek menempuh
        jarak yang sama setiap detik terlepas dari laju bingkainya.
\end{enumerate}

Perangkatnya sudah lengkap. Yang belum ada hanyalah
pemakaiannya. Hal ini disengaja: memisahkan pemasangan siklus dari
penggunaan selang waktu membuat kedua persoalan dapat dipelajari satu per
satu. Kesalahan pada gerak nanti dapat ditelusuri tanpa harus mempertanyakan
apakah siklusnya sudah benar.

\subsection{Langkah Tetap dan Langkah Berubah pada Siklus Ini}

Bagian~2.2 telah membedakan dua cara memajukan simulasi. Sesudah menulis
siklus sendiri, perbedaan itu dapat ditinjau kembali dengan contoh yang nyata.

\textbf{Pada \istilah{fixed timestep}}, simulasi selalu maju dengan panjang
langkah yang sama, misalnya 1/60 detik, sedangkan selisih waktu yang belum
terpakai ditabung pada penampung. Cara ini menghasilkan simulasi yang dapat
diulang persis dan perhitungan fisika yang stabil.

\textbf{Pada \istilah{variable timestep}}, simulasi maju mengikuti lama bingkai
yang sesungguhnya, sehingga panjang langkahnya berubah-ubah. Cara ini lebih
sederhana untuk diterapkan.

Siklus yang baru saja Anda tulis termasuk golongan kedua. Setiap putaran
memakai selang waktu yang benar-benar berlangsung---kadang 16 milidetik,
kadang 34 milidetik---sehingga langkahnya berubah-ubah. Untuk keperluan Blok~I
dan Blok~II, cara ini memadai, dan sebagian besar permainan sederhana memang
memakainya.

Pilihan itu diambil secara sadar, bukan karena ketidaktahuan. Yang perlu Anda
ketahui adalah batas kemampuannya: apabila panjang langkah suatu ketika
menjadi sangat besar, misalnya karena komputer sesaat tersendat, objek yang
bergerak cepat dapat melompat melewati rintangan tanpa terdeteksi. Gejala
itu beserta penanganannya akan Anda jumpai kembali pada Bab~VII, sedangkan
pada Bab~XIII Unity menyediakan langkah tetap untuk perhitungan fisika dan
menyerahkan sisanya pada tahap yang mengikuti lama bingkai.

\begin{catatabox}
Cara memeriksa jenis langkah yang dipakai cukup sederhana: lihat apakah
simulasi memakai selang waktu hasil pengukuran secara langsung. Apabila ya,
itu langkah berubah. Apabila simulasi memakai nilai tetap dan menyimpan
selisih waktunya, itu langkah tetap.
\end{catatabox}

\subsection{Perkembangan Siklus pada Bab-Bab Berikutnya}

Bab-bab selanjutnya akan menambahkan bagian pada siklus yang sama, bukan
menggantinya. Susunan yang akan Anda jumpai adalah sebagai berikut.

\begin{table}[htbp]
\centering
\caption{Bagian yang ditambahkan pada siklus permainan pada bab-bab berikutnya}
\label{tab:perkembangan-siklus}
{\setlength{\tabcolsep}{4pt}\footnotesize%
\begin{tabularx}{\textwidth}{@{}C{1.6cm} L{4.6cm} Y@{}}
\toprule
\headrow \textbf{Bab} & \textbf{Yang ditambahkan} & \textbf{Akibat pada siklus} \\
\midrule
Bab~V & Posisi, kecepatan, dan pemakaian selang waktu &
Tahap pembaruan mulai mengubah keadaan objek, dan penayangan menggambar objek
pada posisinya yang baru. \\
\addlinespace
Bab~VI & Gravitasi, lompatan, dan pantulan &
Tahap pembaruan menghitung percepatan, lalu memperbarui kecepatan sebelum
memperbarui posisi. \\
\addlinespace
Bab~VII & \istilah{Sprite}, \istilah{Group}, dan tabrakan &
Tahap pembaruan memeriksa tabrakan antar-objek dan mengubah keadaan permainan
sebagai akibatnya. \\
\addlinespace
Bab~VIII--XI & Pengendalian NPC &
Tahap pembaruan memanggil lapisan keputusan dan lapisan navigasi, dan sebagian
pekerjaan itu tidak dijalankan setiap bingkai karena alasan kinerja. \\
\addlinespace
Bab~XIII & Fisika mesin pada Unity &
Langkah tetap dijalankan mesin untuk perhitungan fisika, sedangkan kode Anda
dipanggil pada tahap yang mengikuti lama bingkai. \\
\bottomrule
\end{tabularx}}
\end{table}

Dengan gambaran tersebut, dapat dikatakan bahwa bab ini adalah tulang
belakang seluruh Blok~I dan Blok~II. Setiap bab berikutnya mengisi tahap
pembaruan dengan pekerjaan baru, sementara kerangka siklusnya tetap seperti
yang baru saja Anda tulis.
